\section{Entfaltung (Deconvolution) von Messungen}
\label{sec:deconvolution}
Hat man mit einem der in \ref{sec:ir_msrmt} beschriebenen Verfahren das Ausgangssignal eines Raums gemessen, dann ist der nächste wichtige Schritt die Berechnung der
Impulsantwort des Raumes. Laut dem Convolution Theorem \ref{sec:convolution_theorem} ist eine Convolution in der Zeitdomäne äquivalent zu einer Multiplikation in der Frequenzdomäne \cite[Kap. 2]{oppenheim1999}.
Naiv liegt es damit nahe, dass eine Deconvolution in der Zeitdomäne äquivalent zu einer Division in der Frequenzdomäne ist. Mathematisch ist diese Annahme auch korrekt,
allerdings gibt es in der Praxis, besonders in Anwendungen, die später zum Anhören gedacht sind, besser geeignete Verfahren, die mit den potenziell auftretenden nonlinearen
oder zeitvarianten Eigenschaften der gemessenen Systeme besser umgehen könnnen. 